\subsection{An integral Morita descent at multiplicity five}
\label{sec:m5-4982-morita-descent}

The next argument couples the complete right action of a Morita dual,
then descends to a fusion subcategory of dimension $148$. It applies to
arbitrary fusion categories over $\mathbb C$.

\begin{lemma}\label{lem:four-p-cyclic-invertibles}
Let $p>3$ be prime. If an integral fusion category has dimension type
$(1^4,4^{(p-1)/4})$, its group of invertible objects is cyclic of order
four. In particular this dimension type with invertible group
$C_2\times C_2$ cannot occur.
\end{lemma}
\begin{proof}
Its dimension is $4p$. Jordan--Larson's classification
\cite[Theorem~1.1]{JordanLarson} implies that it is group-theoretical: in that theorem
put $q=2$, and note that neither $p=2$ nor $p\mid q+1$ is possible.
Write it as $\mathcal C(G,\omega,H,\psi)$, where $|G|=4p$.
The simples supported on a double coset $HgH$ have dimensions
\[
 [H:H\cap gHg^{-1}]\dim(\pi),
\]
where $\pi$ runs through irreducible projective representations of
$H\cap gHg^{-1}$ for the prescribed cocycle. The simples supported
on $H$ form a fusion subcategory equivalent to $\operatorname{Rep}(H)$;
this statement incorporates the cochain $\psi$ and is independent of
its choice. We use the standard double-coset description
\cite[Theorem~5.1]{GelakiNaidu}.

The Sylow $p$-subgroup $P$ is normal, since its number divides four
and is one modulo $p$. By Schur--Zassenhaus,
$G=P\rtimes K$ with $P\cong C_p$ and $|K|=4$.
If $|H|=p$, the identity double coset already supplies $p>4$
invertibles. If $|H|=2p$, then $H$ is cyclic or dihedral: the former
supplies $2p$ invertibles and the latter supplies dimension-two
simples. Both contradict the given dimension type.
If $H=G$, the category is equivalent to $\operatorname{Rep}(G)$.
Since $K$ is abelian and $P$ is cyclic, the ordinary irreducible
dimensions are the orbit sizes of $K$ on $\widehat P$.
For completeness, each character of $P$ extends to its stabilizer
$P\rtimes K_\chi$ by taking value one on $K_\chi$, so no projective
representation factor changes these sizes. A dimension-four simple
therefore requires a faithful action of $C_4$ on $P$.
Then $[G,G]=P$, and the invertible group
$\widehat{G/[G,G]}$ is $C_4$.

It remains to consider $H$ of order at most four; conjugating $H$
places it in $K$. For $|H|\leq2$, the double-coset formula cannot
produce dimension four. Suppose $|H|=4$.
If $H\cong C_2\times C_2$, its action on $P$ has a kernel $Z$ of
order at least two, because $\operatorname{Aut}(P)$ is cyclic.
This kernel is central in $G$, hence is contained in every
$H\cap gHg^{-1}$. Thus the index is at most two. An intersection of
order two has only one-dimensional projective irreducibles; an
intersection of order four has projective irreducibles of dimension
at most two. Again no simple has dimension four.
If $H\cong C_4$ but its action on $P$ is not faithful, the same kernel
argument applies, and every intersection is cyclic, so all projective
irreducibles are one-dimensional. Finally suppose that $H=C_4$ acts
faithfully. It acts freely on $P\setminus\{1\}$, so
$N_G(H)=H$. A dimension-one simple must have double-coset index one;
hence every invertible lies in the identity-double-coset subcategory
$\operatorname{Rep}(H)$. Its four invertibles form $\widehat H=C_4$,
as required. This argument allows arbitrary $\omega$ and $\psi$.
\end{proof}

\begin{proposition}\label{prop:m5-4982-exclusion}
The fusion ring $R7\_5\_4982$ has no categorification over $\mathbb C$.
\end{proposition}
\begin{proof}
Write its basis as $1,g,g^2,X,Y,Z,W$, where $g^3=1$ and $g$ fixes
the four noninvertibles. These are self-dual, have dimensions
$3,12,12,12$, and, with $S=1+g+g^2$, satisfy
\[
\begin{gathered}
 X^2=S+2X,\qquad XV=3V\quad(V=Y,Z,W),\\
 Y^2=S+3X+2Y+4Z+5W,\quad
 Z^2=S+3X+5Y+2Z+4W,\\
 W^2=S+3X+4Y+5Z+2W,\\
 YZ=4Y+5Z+3W,\quad YW=5Y+3Z+4W,\quad
 ZW=3Y+4Z+5W.
\end{gathered}
\]
Suppose that a categorification $\mathcal C$ exists.
The action of the pointed subcategory on the single simple $X$
trivializes its associator class, so
$A=1\oplus g\oplus g^2$ is a connected separable algebra.
There are thirteen simple right $A$-modules: $A$ itself, and three
lifts of each of $X,Y,Z,W$. Their normalized dimension vector is
\[
 d=(1,1,1,1,4,4,4,4,4,4,4,4,4).
\]
The algebra automorphism twisting by a primitive cube root
$\zeta$ cyclically permutes each set of three lifts and fixes $A$.
Use labels $0$ for $A$, $1,2,3$ for the $X$ lifts, and consecutive
triples starting at $4,7,10$ for the other lifts.

\emph{Exhausting the module actions.}
On a nontrivial Fourier sector of the twisting action, a source
simple has a $4\times4$ matrix over $\mathbb Z[\zeta]$.
If the total multiplicity between two underlying fixed simples is
$n$, each entry is of the form
\begin{equation}
 a+b\zeta+c\zeta^2,\qquad a,b,c\in\mathbb Z_{\geq0},\quad a+b+c=n.
 \label{eq:4982-Fourier-domain}
\end{equation}
Duality makes the four noninvertible matrices Hermitian. The matrix
of $g$ is diagonal with cube-root entries. If a diagonal entry were
not one, the relation $gV=V$ would force the corresponding row of
every noninvertible matrix to vanish. This is impossible: for the
$X$ coordinate use $N_{XX}^X=2$, and for the other three coordinates
use the diagonal entries $2,5,4$ of multiplication by $Y$. None is
divisible by three, whereas a zero entry in \eqref{eq:4982-Fourier-domain} requires that
$n$ be divisible by three. Thus $T_g=I$.
The fixed-simple block of multiplication by $X$ is
$\operatorname{diag}(2,3,3,3)$. Its off-diagonal entries vanish, and
$T_X^2=3I+2T_X$. The real domains in \eqref{eq:4982-Fourier-domain} therefore force
\[
 T_X=\operatorname{diag}(-1,3,3,3).
\]
Commutation splits $T_Y=0\oplus T$, where $T$ is Hermitian of size
three and its entries satisfy \eqref{eq:4982-Fourier-domain} with majorant
\[
 B=\begin{pmatrix}2&4&5\\4&5&3\\5&3&4\end{pmatrix}.
\]
The raw fusion tensor gives the annihilator
\[
 T(T-12I)(T^3+T^2-12T+11I)=0.
\]
The leading principal minors of $12I-B$ are $10,54,47$, so $B$
has spectral radius less than twelve. The entrywise bound
$|T_{ab}|\leq B_{ab}$ excludes eigenvalue twelve.
The characteristic polynomial of $T$ has rational integral
coefficients: it is real, has coefficients in $\mathbb Q(\zeta)$,
and all its eigenvalues are algebraic integers. The cubic
$t^3+t^2-12t+11$ is irreducible. Consequently either the entire
three-dimensional spectrum is zero, or its characteristic polynomial
is this cubic. The zero case contradicts
$\operatorname{tr}(T)\equiv\operatorname{tr}(B)=11\pmod3$.
The necessary moments are therefore
\[
 \operatorname{tr}(T)=-1,\qquad
 \operatorname{tr}(T^2)=25,\qquad
 \operatorname{tr}(T^3)=-70.
\]
Exact enumeration of the finite entry domains \eqref{eq:4982-Fourier-domain} gives $639$
matrices for the first two moments and eighteen for all three.
On this sector the remaining action matrices are forced by
\[
 T_Z=7I-3T-T^2,\qquad T_W=-8I+2T+T^2;
\]
the identities are verified in the cubic spectral summand of the
original regular representation. Each of the eighteen choices
reconstructs nonnegative integral $13\times13$ module matrices
$M_i$ satisfying every source fusion relation. Explicit basis
permutations fixing label zero identify all eighteen module actions.
In particular it suffices to use the first one. It satisfies
\[
 M_g=M_{g^2}=I_{13},\qquad
 M_X=(J_4-I_4)\oplus3I_9.
\]

\emph{Exhausting the complete Morita right action.}
Let $\mathcal D={}_A\mathcal C_A$, and let $B_U$ denote the right
module action of a simple $U$ of $\mathcal D$.
Associativity in the Morita bicategory gives the integral Gram
identity
\begin{equation}
 G_{(a,b),(c,e)}=\sum_i(M_i)_{c,a}(M_i)_{e,b}
   =\sum_U(B_U)_{a,b}(B_U)_{c,e}.
 \label{eq:4982-full-Gram}
\end{equation}
Indeed both sides compute the same multiplicity by associating
$m_a(m_c^*m_e)$ in the two possible ways. Equivalently, $G$ is the
Gram matrix of the vectorized right actions. They lie in
$\operatorname{im}_{\mathbb Q}G$, commute with the $M_i$, are
permuted by transposition, and satisfy the entrywise bounds
\begin{equation}
 0\leq (B_U)_{a,b}\leq
 \left\lfloor\sqrt{G_{(a,b),(a,b)}}\right\rfloor,
 \qquad \operatorname{vec}(B_U)\operatorname{vec}(B_U)^t\leq G.
 \label{eq:4982-action-bounds}
\end{equation}
Here $\operatorname{rank}(G)=37$.
The restriction of $G$ to pairs $(a,0)$ is $3I_{13}$.
Every first column of $B_U$ is therefore a unit vector, and each
unit vector occurs three times. No first column is zero: label
zero generates the indecomposable module category and a commuting
nonzero right action cannot annihilate it. If the column is $e_a$,
then $\operatorname{FPdim}(U)=d_a$. Hence $\mathcal D$ has dimension
type $(1^{12},4^{27})$. Transposition gives the same statement for
first rows.

For a noninvertible action, write its first column and row as
$e_a$ and $e_b^t$, with $a,b\geq4$. Since
$B_U M_X e_0=M_X B_U e_0=3e_a$, its first four by four block is
zero, by nonnegativity. The cross-Gram block
\begin{equation}
 \bigl(G_{(a,0),(0,b)}\bigr)_{a,b=4}^{12}
      =\operatorname{diag}(J_3,J_3,J_3)
 \label{eq:4982-cross-Gram}
\end{equation}
forces $a,b$ to lie in the same triple and forces exactly one
noninvertible simple for every such ordered pair.
The permutation centralizer of the $M_i$ has order $72$; its
stabilizer of label zero has order $18$ and is transitive on each
of these three triples. Independently multiplying an action on the
left and right by elements of this stabilizer preserves \eqref{eq:4982-full-Gram}--\eqref{eq:4982-action-bounds},
and moves its first column and row to a common anchor $4$, $7$, or
$10$. These multiplications are used only to enumerate necessary
matrices, without asserting that the permutations are tensor functors.

For each anchor, the affine subspace in
$\operatorname{im}_{\mathbb Q}G$ with that first row, first column,
and zero small block has dimension twelve. Selecting twelve actual
matrix entries as coordinates gives a common denominator $31$.
Integrality imposes two independent congruences modulo $31$; their
solved coordinates have bounds at most four, so each residue has
at most one representative. Each of the three remaining integer
boxes has exactly $2{,}048{,}000$ points. Checking \eqref{eq:4982-action-bounds} retains
exactly one matrix per anchor. Stabilizer translates give exactly
$27$ distinct matrices, one per allowed ordered pair in \eqref{eq:4982-cross-Gram}, so all
occur once.

Subtract their outer products from $G$. Every invertible right
action is a permutation commuting with the $M_i$, hence lies in
the above group of order $72$. Its image is a subgroup $H$ of order
dividing twelve; a possible kernel would make every element of $H$
occur with the same multiplicity $12/|H|$. Exhaustive subgroup
closure gives respectively $1,9,13,7,13,10$ subgroups of orders
$1,2,3,4,6,12$. Exactly two satisfy the remaining Gram identity,
and both have order twelve. Thus the invertible right actions are
faithful in either case. This step does not assume that all $39$
action matrices are linearly independent; their span has dimension
only $37$.

\emph{Recovering the forced fusion subcategory.}
For each of the two complete frames, products involving an
invertible are read directly from matrix multiplication. For two
noninvertibles $U,V$, the invertible coefficient of $h$ is forced
by Frobenius reciprocity:
\[
 [UV:h]=[U^*h:V].
\]
Subtract these twelve coefficients from $B_UB_V$. Its first
column gives, for each of the nine large module coordinates, the
sum of the remaining three coefficients with that first column.
Their total is at most four, by the dimension equation
$16=\sum_h[UV:h]+4\sum_W[UV:W]$.
Enumerating these nonnegative compositions and checking all matrix
entries gives exactly one solution for each of the $27^2=729$
products, in both frames. The resulting two rank-$39$ tensors
satisfy the based-ring axioms, including associativity and both
Frobenius reciprocity identities.

In the recorded ordering of each frame, the simple with index $14$
generates the same closed subset
\[
 \{0,1,4,9,14,15,18,23,25,27,30,33,38\}.
\]
Its dimension type is $(1^4,4^9)$ and its dimension is $148$.
The four invertibles have indices $0,1,4,9$, and all are
self-inverse, so their group is $C_2\times C_2$.
A based subset closed under products and duals gives a fusion
subcategory of a hypothetical categorification. Applying
Lemma~\ref{lem:four-p-cyclic-invertibles} with $p=37$ gives the
contradiction.

Every finite enumeration in this proof is regenerated from the
original seven-dimensional tensor by
\path{evidence/rank7_m5_20261020/modules/verify_r4982.py}.
The exhaustive Fourier data, normalized bounded domains, two
right-action frames, fusion tensors, and closed-subring tensors
are retained with the final certificate
\texttt{r4982\_exclusion\_certificate.json}. The verifier uses exact
integer or rational arithmetic throughout; it invokes neither a
numerical feasibility test nor a pivotality assumption.
\end{proof}
